% 阶段报告三导读。编译：xelatex summary_plain_zh.tex（两遍）。
\documentclass[12pt,a4paper]{article}
\linespread{1.10}
\usepackage[margin=0.9in]{geometry}
\usepackage{amsmath,booktabs,tabularx,graphicx,xeCJK}
\setCJKmainfont{FandolSong-Regular}[Extension=.otf,BoldFont=FandolSong-Bold]
\setCJKsansfont{FandolHei-Regular}[Extension=.otf,BoldFont=FandolHei-Bold]
\usepackage[hidelinks]{hyperref}
\renewcommand{\figurename}{图}
\renewcommand{\tablename}{表}
\emergencystretch=3em
\setlength{\parskip}{4pt}
\title{\bfseries 补回漏掉的细胞核，怎样少添一些错误？\\[7pt]
\large 阶段报告三导读}
\author{课程项目 bdd26}
\date{2026 年 10 月 6 日修订}
\begin{document}
\maketitle

\section*{先说目前的结论}

我们找到了一个小幅改进的办法：保留原尺度模型的分割结果，用高尺度模型补入可能漏掉的核，再过滤掉其中一部分特别小的候选。这样能保留大部分 Dead 类核的得分提升，同时让整体分割得分的平均损失接近零。

这里有两个限定。第一，\textbf{平均损失很小，不等于已经证明没有损失}。第二，方法需要运行两种尺度的模型，还没有测过完整流程要多花多少时间。因此，目前适合把它作为一个有依据的候选方案，不能说已经全面胜过基线。

\section*{为什么要同时使用两种尺度}

这个项目让模型在病理图像中逐个圈出细胞核，再判断它们的类别。Dead 是数据集中的一类，主要对应死亡、凋亡及相关碎裂核形态。这类核较少，也较难识别。

前两阶段发现，把同一张图像从 $256\times256$ 插值到 $512\times512$，再配套训练模型，确实有助于找到一些原来漏掉的核。但它也会引入额外错误，导致整体分割成绩下降。放大并没有增加拍摄时不存在的细节，它只是改变了模型处理图像的尺度。

于是，本阶段尝试保留两者各自有用的部分：以 CellViT-UNI 为基础，保留原尺度模型经过后处理的结果；只从放大尺度模型里挑出不与现有核重叠的候选，补到空白处。这个方法记为 P2。

不过，在开展新实验前，我们先纠正了旧结果。旧后处理有实现错误，原来的九份训练结果也有种子重复，不能当成完整的三种子检验。修复后，我们补齐了缺失训练，分别在三个数据划分、三个随机种子下重复训练和比较同一套流程；随机种子用于控制初始化等训练随机过程。九次比较是九个模型配对，并不是九批独立病例。

\section*{第一轮结果：补得回来，但错误仍然偏多}

P2 用面积、类型置信度和是否贴边来筛选候选，规则只在验证集上选择；没有合适规则时，就什么也不添加。匹配种子实验中，Dead 得分八对提高、一对不变，但整体分割指标 bPQ 的平均损失超过了事先约定的范围，所以这一轮没有通过。

此前，我们对随机种子 19 下两个有新增候选的划分（划分 2 和 3）做过诊断。这两项测试集统计提示了一个可能原因：
\begin{itemize}
\item 新增候选中，只有约 \textbf{39--43\%} 能与人工标注的某个核较好地重合；
\item 在这些已经匹配上的候选中，约 \textbf{86--88\%} 的类型判断正确。
\end{itemize}

第二个数字并不差，但不能拿它代替第一个数字。模型很有把握地判断一个候选属于哪类，不代表它画出的那个对象本身就是一个合格的核实例。未匹配候选可能是误检、碎片，也可能是轮廓偏差过大；它们都会在评测中造成损失。

\section*{第二轮只改一件事：过滤特别小的添加}

接下来我们只检查验证集，不再从测试图里寻找新阈值。九个匹配对一共添加了 17,359 个候选，其中 41.50\% 能匹配人工标注。按面积分组，区别很明显：

\begin{center}
\begin{tabular}{lr}\toprule
新增候选的面积 & 与标注匹配的比例\\\midrule
不足 30 个像素 & 19.8\%\\
30--59 个像素 & 44.4\%\\
60--99 个像素 & 48.9\%\\\bottomrule
\end{tabular}
\end{center}

这里说的是原始 $256\times256$ 网格中的\textbf{面积}，不是边长。最小的一组约八成没有匹配上，过滤它们有望少添错误。但其中也有真核，包括 43 个正确匹配并分为 Dead 的候选，所以面积过滤也有误删正确候选的风险。

我们在实验前固定了五种选择：不添加；保持原 P2；添加面积至少 30；非 Dead 至少 30、Dead 例外；添加面积至少 60。仍由各自的验证集选择，不按测试成绩挑规则。最后，七对选了面积至少 30，一对给 Dead 留了例外，另一对继续不添加。加上这个步骤的方法记为 EF-P2。

\section*{改善有多大}

下面都是三个划分各自平均三个种子、再对划分取平均的测试得分，按 0--1 计。基线是原尺度模型经过相同验证集调参流程选择的后处理结果，并非未经处理的原始预测。

\begin{center}
\begin{tabular}{lrrr}\toprule
方法 & 综合指标 mPQ & 整体分割 bPQ & Dead PQ\\\midrule
原尺度公平基线 & 0.49760 & 0.66530 & 0.17255\\
P2 添加 & 0.49769 & 0.66443 & 0.17997\\
EF-P2 添加并过滤 & 0.49828 & 0.66529 & 0.17954\\\bottomrule
\end{tabular}
\end{center}

过滤后，Dead 得分仍比基线高 0.00699，保留了原 P2 约 \textbf{94\%} 的提升。bPQ 与基线的平均差距则从约 $-0.00087$ 缩小到 $-0.000009$。在八对确实做了添加的比较中，过滤后的 mPQ 和 bPQ 都比未过滤 P2 高。

EF-P2 通过了事先约定的判定规则，成为当前保留的候选方案。过滤后，Dead PQ 仍保留原有的大部分增益；这是得分变化，不能直接理解成保留了相同比例的正确实例。相对基线，Dead 得分八对提高、一对不变。

\section*{为什么还不能说问题解决了}

一是改善并不均匀。Dead 的提升在第三个划分最明显，第二个划分仍有较大不确定性。整体分割和综合指标的区间也都跨过零，即现有数据仍容许改善或下降两种可能，不能排除一定程度的损失。把各类在缺少相应真值的图像上的误报也计入后，严格总体指标仍略有下降。

二是同一批测试图已经参与了多轮研究。每一轮都先固定规则、只用验证集选参，是必要的，但不能把多次看过的测试数据重新称为独立验证。多个模型重复处理同一张图，也不会增加独立病例的数量。

三是尚未测量成本。面积过滤本身简单，但候选来自另一套高尺度模型。最终是否值得使用，还要看额外推理时间、资源占用以及新数据上的表现。

我们还纠正了一项事后的边界统计。重新计算发现，贴边候选的总体匹配率反而高于内部候选；但正确补回的 Dead 主要在内部。这提醒我们，不能因为对象贴边就一概删除，也不能把某个分组的平均差异当作因果规律。这个事后发现没有用于更改本轮规则。

\section*{下一步最值得做什么}

现在不宜继续围着同一批测试图增加阈值。更有价值的是固定 EF-P2，在新的独立数据或严格重新设计的留出实验中检验，并同时测量两套模型的完整运行成本。若暂时没有独立数据，就把目前结果保留为阶段性成果，不把小幅改善包装成定论。

本阶段把一个笼统问题变得具体了：不仅要问模型能否多找到核，还要问\textbf{新增的候选有多少值得保留}。把“是什么类别”和“是不是可靠实例”分开检查，正是这次改进的关键。

\vspace{1em}
\noindent\small 完整数值、置信区间、历史更正和复现说明见同目录下的
\texttt{stage\_report\_3\_zh.pdf}；英文版为 \texttt{stage\_report\_3.pdf}。
\end{document}
